% ID: MAT-ALG-PER-008
% Titolo: Percentuali inverse e variazioni: sei sfide
% Disciplina: Matematica
% Area: algebra
% Argomento: percentuali
% Classe: Prima secondaria di secondo grado
% Difficolta: 3
% Tipo: laboratorio
% Tag: matematica, algebra, percentuali
% Fonte: laboratorio_percentuali_anonimizzato
% Provenienza: materiale fornito dal docente; testo anonimizzato; soluzione redatta per Exergo

\begin{esercizio}
Risolvi, anche con un foglio di calcolo:
\begin{enumerate}
\item Quanto risparmi con il $15\%$ di sconto su 48 euro?
\item Qual è il prezzo finale dello stesso oggetto?
\item Quale percentuale rappresentano 12 figurine rare su 40?
\item Un prezzo passa da 80 a 100 euro: qual è l'aumento percentuale?
\item 100 euro aumentano del $20\%$ e poi diminuiscono del $20\%$: quanto rimane?
\item Una felpa costa 54 euro dopo lo sconto del $10\%$: qual era il prezzo iniziale?
\end{enumerate}
\end{esercizio}

\begin{soluzione}
I calcoli sono, nell'ordine: $48\cdot0{,}15=7{,}20$ euro;
$48\cdot0{,}85=40{,}80$ euro; $(12/40)100\%=30\%$;
$[(100-80)/80]100\%=25\%$; $100\cdot1{,}2\cdot0{,}8=96$ euro;
$54/0{,}9=60$ euro. Nel problema inverso non si aggiunge il $10\%$ a 54:
la base dello sconto è il prezzo originale.
\end{soluzione}
